% \iffalse
%<*driver>
\documentclass{l3doc}
\usepackage[brazilian]{babel}
\usepackage{booktabs}
\begin{document}
  \sloppy
  \DocInput{abntex3-manual.dtx}
\end{document}
%</driver>
% \fi
%
% \title{Manual do \pkg{abntex3} Community}
% \author{Rafael dos Santos Dantas}
% \date{2026-08-03 v0.1.0-alpha.1}
%
% \maketitle
% \tableofcontents
%
% \begin{documentation}
%
% \section{Apresentação}
%
% O \pkg{abntex3} Community é um pacote modular para produzir documentos
% brasileiros sobre as classes usuais \cls{article}, \cls{report} e
% \cls{book}. Esta primeira alfa é destinada a testes: não é recomendada para
% produção, pode receber mudanças de API e não constitui alegação de
% conformidade com normas da ABNT. LuaLaTeX é o único motor suportado.
%
% O pacote cuida da estrutura e apresentação documental. Citações e
% referências são responsabilidade do \pkg{biblatex-abnt}, processado por
% \texttt{biber}. Regras institucionais específicas devem permanecer no
% documento ou em extensões externas ao núcleo.
%
% \section{Instalação}
%
% Para testar uma release, baixe \file{abntex3.tds.zip} e
% \file{SHA256SUMS} na página de Releases, confira o resumo SHA-256 e extraia o
% TDS em uma árvore TeX de ensaio. Também é possível instalar um checkout:
%
% \begin{verbatim}
% git clone https://github.com/dantasrs/abntex3-community.git
% cd abntex3-community
% l3build install
% \end{verbatim}
%
% A distribuição TeX precisa oferecer LuaLaTeX. Documentos com bibliografia
% também precisam de \pkg{biblatex-abnt} e \texttt{biber}. O pacote ainda não
% está no CTAN; os artefatos alfa não devem substituir silenciosamente uma
% instalação usada em produção.
%
% \section{Primeiro documento}
%
% \begin{verbatim}
% \documentclass{report}
% \usepackage[brazilian]{babel}
% \usepackage[perfil=academico,tipo-trabalho=dissertacao]{abntex3}
%
% \ABNTEXsetup{
%   titulo={Documento de ensaio},
%   autoria={Nome da autoria},
%   instituicao={Instituição},
%   natureza={Dissertação},
%   objetivo={apresentada para obtenção do título},
%   area={Área de concentração},
%   orientacao={Nome da orientação},
%   local={Cidade},
%   data={2026}
% }
%
% \begin{document}
% \ABNTEXacademicoexterna
% \ABNTEXacademicopretextual
% \ABNTEXfolhaderosto
% \ABNTEXsumario
% \ABNTEXacademicotextual
% \chapter{Introdução}
% Texto do documento.
% \end{document}
% \end{verbatim}
%
% As opções de carregamento são processadas primeiro. Chamadas posteriores de
% \cs{ABNTEXsetup} têm precedência; em cada nível, a última atribuição de uma
% chave prevalece.
% \cs{ABNTEXvalidate} verifica os metadados técnicos do perfil selecionado.
%
% \section{Classe, perfil e tipo}
%
% A classe fornece a estrutura básica do LaTeX. O perfil coordena o fluxo e os
% elementos de um gênero documental. O tipo é uma variação dentro de um
% perfil. As combinações usuais são:
%
% \begin{center}
% \begin{tabular}{lll}
% \toprule
% Documento & Classe sugerida & Configuração \\
% \midrule
% tese, dissertação, TCC & \cls{report} ou \cls{book} &
%   \texttt{perfil=academico} \\
% projeto de pesquisa & \cls{report} ou \cls{book} &
%   \texttt{perfil=projeto} \\
% artigo & \cls{article} & \texttt{perfil=artigo} \\
% relatório & \cls{report} ou \cls{book} &
%   \texttt{perfil=relatorio} \\
% livro ou folheto & \cls{book} ou \cls{report} &
%   \texttt{perfil=livro} \\
% \bottomrule
% \end{tabular}
% \end{center}
%
% O perfil \texttt{nenhum} expõe módulos gerais sem impor um fluxo. As classes
% \cls{abntex2}, \cls{memoir} e KOMA-Script não são suportadas nesta alfa.
%
% \section{Configuração comum}
%
% \cs{ABNTEXsetup}\marg{chaves} reúne metadados, layout e escolhas comuns.
% Entre as chaves centrais estão \texttt{titulo}, \texttt{subtitulo},
% \texttt{autoria}, \texttt{instituicao}, \texttt{natureza},
% \texttt{objetivo}, \texttt{area}, \texttt{orientacao},
% \texttt{coorientacao}, \texttt{volume}, \texttt{local} e \texttt{data}.
% \cs{ABNTEXmetadata}\marg{nome} recupera um metadado documentado;
% \cs{ABNTEXprofile} e \cs{ABNTEXclass} informam perfil e classe detectados.
%
% Para o PDF, use \texttt{assunto-pdf}, \texttt{palavras-chave-pdf} e
% \texttt{idioma-pdf}. Com \pkg{hyperref} carregado, o pacote transfere título,
% autoria e esses valores aos metadados do arquivo.
%
% \section{Layout, estrutura e fases}
%
% As chaves \texttt{papel}, \texttt{orientacao-papel}, \texttt{impressao} e
% as quatro margens controlam a página. \cs{ABNTEXpretextualpages} inicia a
% contagem não exibida; \cs{ABNTEXtextualpages} passa a exibir a paginação sem
% reiniciar o contador. O ambiente \env{ABNTEXcitacaolonga},
% \cs{ABNTEXfonte} e \cs{ABNTEXnota} tratam exceções tipográficas comuns.
%
% \cs{ABNTEXstructureSetup} configura níveis de seção.
% \cs{ABNTEXsumario} produz o sumário, e \cs{ABNTEXappendix} e
% \cs{ABNTEXannex} iniciam apêndices e anexos. Títulos pré e pós-textuais usam
% \cs{ABNTEXpretextualtitle} e \cs{ABNTEXposttextualtitle}.
%
% \section{Elementos pré-textuais}
%
% O módulo geral oferece:
% \begin{itemize}
% \item \cs{ABNTEXcapa}, \cs{ABNTEXlombada} e \cs{ABNTEXfolhaderosto};
% \item \cs{ABNTEXfichacatalografica}, \cs{ABNTEXerrata} e
%   \cs{ABNTEXfolhadeaprovacao};
% \item \cs{ABNTEXdedicatoria}, \cs{ABNTEXagradecimentos} e
%   \cs{ABNTEXepigrafe};
% \item \cs{ABNTEXresumo} e \cs{ABNTEXresumoemlinguaestrangeira}.
% \end{itemize}
% Dados profissionais de catalogação devem ser preparados externamente.
%
% Há listas automáticas de ilustrações e tabelas e registros declarativos para
% abreviaturas, siglas e símbolos. \cs{ABNTEXvalidatepretextual} diagnostica
% ordem e omissões observáveis no perfil acadêmico.
%
% \section{Bibliografia e elementos pós-textuais}
%
% Escolha \texttt{citacoes=externo}, \texttt{autor-data} ou
% \texttt{numerico}. Os dois últimos carregam \pkg{biblatex-abnt}; o primeiro
% deixa a configuração inteiramente com o documento. Um fluxo típico é:
%
% \begin{verbatim}
% \usepackage[citacoes=autor-data]{abntex3}
% \addbibresource{referencias.bib}
% ...
% \ABNTEXposttextualpages
% \ABNTEXbibliografia
% \end{verbatim}
%
% Compile documentos bibliográficos com LuaLaTeX, \texttt{biber} e mais duas
% passagens de LuaLaTeX. \cs{ABNTEXglossario} e \cs{ABNTEXindice} são pontos
% de integração para os geradores correspondentes.
%
% \section{Perfis documentais}
%
% Cada perfil possui comandos para suas fases, validação e dados próprios:
%
% \begin{description}
% \item[Acadêmico] Comandos de fases com prefixo
%   \texttt{ABNTEXacademico} e \cs{ABNTEXvalidateacademico}.
% \item[Projeto] Comandos com prefixo \texttt{ABNTEXprojeto}, incluindo as
%   cinco divisões textuais e \cs{ABNTEXvalidateprojeto}.
% \item[Artigo] \cs{ABNTEXartigoSetup}, autorias repetíveis, fases, divisões,
%   uma ou duas colunas e validação própria.
% \item[Relatório] \cs{ABNTEXrelatorioSetup}, equipe repetível, elementos
%   próprios, fases e validação própria.
% \item[Livro] \cs{ABNTEXlivroSetup}, autorias e colaborações repetíveis,
%   elementos editoriais, divisões e validação própria.
% \end{description}
%
% Os manuais técnicos de cada perfil descrevem argumentos, obrigatoriedade e
% condicionais. Os modelos canônicos são a referência prática mais completa.
%
% \section{Migração do abnTeX2}
%
% O \pkg{abntex3} é um pacote, não uma classe. Migre primeiro a classe-base e
% os metadados, depois as fases e elementos, e por fim a bibliografia para
% \pkg{biblatex-abnt}/\texttt{biber}. O módulo opcional
% \pkg{abntex3-compat} oferece apenas aliases inequívocos e emite avisos; ele
% não pretende reproduzir \cls{memoir}, \pkg{abntex2cite} ou toda a superfície
% histórica do abnTeX2.
%
% \section{Modelos, testes e diagnóstico}
%
% O diretório \file{examples/} contém modelos completos para todos os perfis.
% Comece por uma cópia do modelo correspondente e remova apenas os elementos
% opcionais de que não precisa. Para diagnosticar uma configuração, ative
% \texttt{debug=true} e execute os validadores do perfil.
%
% Desenvolvedores devem rodar pelo menos:
%
% \begin{verbatim}
% git diff --check
% l3build check
% l3build doc
% \end{verbatim}
%
% Correções observáveis precisam de teste que falhe antes da correção e passe
% depois. Alterações de distribuição também exigem \texttt{l3build ctan} e
% \texttt{scripts/validate-distribution.sh}; mudanças no PDF canônico exigem
% \texttt{scripts/quality-audit.sh}.
%
% \section{Estabilidade e como colaborar}
%
% A baseline da primeira alfa protege os nomes públicos contra remoção
% acidental, mas a linha \texttt{0.x} ainda pode mudar. Alterações
% incompatíveis exigem discussão, changelog, migração, teste e nova versão.
% Relatos de uso devem informar sistema, TeX Live, classe, perfil e exemplo
% mínimo sem dados sensíveis e mencionar a issue 23. Divergências normativas
% devem citar norma, edição e item sem reproduzir conteúdo protegido; a
% revisão independente das matrizes é acompanhada na issue 24.
%
% A referência completa, o guia alfa, a política de API, as matrizes de
% requisitos e os canais de contribuição estão no repositório:
% \url{https://github.com/dantasrs/abntex3-community}.
%
% \end{documentation}
%
% \Finale
